This section introduces Vision Transformers (ViTs) and the conceptual machinery used throughout the survey, covering patch tokenization, multi-head self-attention, and the ViT-versus-CNN comparison. Foundational methods include ViT (2020, patch tokenization plus standard Transformer encoder), DeiT (2020, distillation-token data-efficient training), T2T-ViT (2021, progressive token-to-token re-tokenization), Swin (2021, shifted-window hierarchical attention), CaiT (2021, LayerScale and class-attention), MAE (2021, 75\%-mask pixel reconstruction), DINO (2021, EMA-teacher self-distillation), CLIP (2021, image-text contrastive pretraining), Mask2Former (2022, masked-attention universal segmentation), DINOv2 (2023, curated 142M-image self-distillation), ViT-22B (2023, 21.7B-parameter scaling), SAM (2023, promptable segmentation on SA-1B), SAM 2 (2024, video promptable segmentation with memory), and ViT-5 (2026, modernized canonical recipe). Vision Transformers (ViTs) are deep neural networks that treat an image as a sequence of small patches. They process that sequence with the same Transformer encoder originally designed for machine translation. The original Vision Transformer was introduced by Dosovitskiy et al.~in the ICLR 2021 paper ``An Image is Worth 16×16 Words''. It showed that a pure Transformer with neither convolutions nor explicit two-dimensional locality priors can match or exceed the best convolutional networks of the time. The decisive condition is sufficient pretra\ldots{}
